\section{Actual lattice and adelic blocks for the original global setting}
\label{sec:adelic-setting-v47}

This chapter constructs the original setting blocks, rather than assuming
their existence as part of a completed \texttt{SG.HidaSetting}.
Throughout, \(D\) is the original \(D7\) record, \(E/F\) its CM
extension, \(V=\mathrm{VS}(D)\), \(\Lambda=V.L_0\),
\(s=D.R.rS\), \(k=D.R.k\), and \(n=2s+k+1\).
The Hermitian form is conjugate-linear in its first variable.
We retain the original \(\mathrm{Forms0}(D,D.u_0)\) precisely where
the lattice producer requires it, and the original \(\mathrm{GapHyp}\)
precisely where the discrete-series bundle requires it.
Neither condition is replaced by a new existence hypothesis.
The source is pinned at
\texttt{7cbf7a0c32fca1badf05931879a37c3f8a2b03b4}.

% BEGIN ASG1
\paragraph{ASG1. The actual unitary-group instance on the original vector space.}
\label{claim:ASG1}
Every original \(D\) has a \(\mathrm{UnitaryGroupInstance}(D)\)
with coefficient fields exactly \(E,F\), and with the original lattice
and form in the original frame coordinates.
\begin{proof}
Let \(b_i\) be \(V.\mathrm{frame}\), put
\(J_{ij}=h(b_i,b_j)\), and let
\(c:E^n\longrightarrow V\) send \(x\) to \(\sum_i x_i b_i\).
Set \(L=c^{-1}\Lambda\).  Expanding both variables gives
\[
                 \overline{x}^{T}Jy=h(c(x),c(y)).
\tag{ASG.1}
\]
Thus \(J\) is Hermitian.  If \(Jx=0\), (ASG.1) and
nondegeneracy give \(c(x)=0\), hence \(x=0\); consequently
\(\det J\ne0\).  Finite generation and full span of \(L\)
follow by transporting those properties of \(\Lambda\) through
\(c\).  The original unimodularity assertion becomes exactly
\[
 x\in L\quad\Longleftrightarrow\quad
       h_J(y,x)\in\mathcal O_E\text{ for every }y\in L.
\tag{ASG.2}
\]

Use the restricted-product topology on the finite adelic unitary group
in these coordinates.  Local integral stabilizers are compact open:
after choosing local lattice bases they are closed in integral general
linear groups and contain a congruence neighbourhood of the identity.
The resulting congruence subgroups give the exact topology fields of
\(\mathrm{AutInst}\), as in FV4.

The CM pairs of embeddings are in bijection with real places of \(F\).
Use the original chosen member \(\mathrm{embS}(D,\tau)\) of each
pair.  Its restriction is the corresponding real embedding.  The
original \(D.sp.Y\_sig\) field gives signature
\((k+s,s+1)\).  Hermitian diagonalization therefore gives a matrix
\(P_v\) with
\(P_v^*\phi_v(J)P_v=J_{\epsilon_v}\).
The original Boolean sequence \(\mathrm{epsS}\) has exactly
\(k+s\) positive and \(s+1\) negative entries; permuting a
diagonalizing basis puts the signs in that literal order.
This supplies \(\epsilon,\phi,P\) and their identities.

Take the coefficient isomorphism \(\iota_E\) to be the identity,
the basis to be \((b_i)\), and the place bijection to be the inverse
of the original \(\mathrm{tauPl}\).  Complex conjugation agrees
with the CM involution, and its fixed field is \(F\).
Equation (ASG.1), the definition of \(L\), the rank identity in
\(D\), and the preceding sign count verify every remaining field
of \(\mathrm{UnitaryGroupInstance}\).  No lattice is replaced
by a principal one.
\end{proof}
% END ASG1

% BEGIN ASG2
\paragraph{ASG2. The original trace-one element and prescribed fractional ideals.}
\label{claim:ASG2}
There is \(t\in\mathcal O_E\) with \(t+\bar t=1\).
Every prescribed ideal
\(\mathfrak a_i=\mathrm{aaD}(D,i)\) is nonzero and invariant
under conjugation.  At every prime above
\(g=D.host.g(0)D.host.g(1)\), both \(\mathfrak a_i\) and its
inverse are the unit local ideal.
\begin{proof}
The original \(\mathrm{cmUnram}\) field gives, at every maximal
ideal of \(\mathcal O_E\), an integral \(a\) for which
\(a-\bar a\) is outside that ideal.  At a ramified quadratic
place, conjugation would lie in inertia and fix every residue, which
contradicts this field.  Thus \(E/F\) is unramified at finite
places.  The trace on the corresponding integral local extension is
surjective: in the split case it is addition on two copies of the local
ring; in the unramified field case its reduction is the surjective
finite-field trace, and a lift of an element with nonzero trace has
unit trace.  The integral trace image is an ideal of \(\mathcal O_F\)
which is the unit ideal at every maximal ideal, hence is the unit ideal.
Choose \(t\) of trace one.  This argument includes residue
characteristic two.

The ideals \(\mathfrak a_i\) are the specified integral powers,
possibly negative, of the different of \(F/\mathbb Q\) extended
to \(E\).  The different is nonzero and comes from the fixed
field, so it and all its powers are conjugation invariant.
Each of the two original rational primes splits completely in \(E\):
the source counts \([E:\mathbb Q]\) distinct primes above it,
and the identity \(\sum e_w f_w=[E:\mathbb Q]\) forces every
ramification and residue degree to be one.  The primes are therefore
unramified in \(F\), so its different is a unit at all these
places.  This proves the final assertion without excluding either
original prime.
\end{proof}
% END ASG2

% BEGIN ASG3
\paragraph{ASG3. A split prime ideal in the required ideal class.}
\label{claim:ASG3}
For every nonzero fractional ideal \(\mathfrak c\) of a CM
field \(E\), and every nonzero \(N\in\mathcal O_E\), there
are an integral prime ideal \(\mathfrak p\) and \(\alpha\ne0\)
such that
\[
 \mathfrak p+\overline{\mathfrak p}=\mathcal O_E,
 \qquad \mathfrak p+(N)=\mathcal O_E,
 \qquad \mathfrak p\mathfrak c=(\alpha).
\tag{ASG.3}
\]
\emph{Classical input and proof of the exact application.}
Use the Hilbert class field and Chebotarev's theorem in their usual
unconditional forms.  A primary exposition is J.~S.~Milne,
\emph{Class Field Theory}, version 4.03, Chapter V, \S3
(especially Example 3.9), and Chapter VIII, Theorem 7.4:
\url{https://www.jmilne.org/math/CourseNotes/CFT.pdf}.
The two inputs say that the ideal-class Artin map identifies the
class group with the unramified abelian Galois group, and that a
conjugacy class in a finite Galois group occurs at infinitely many
unramified primes.  They are named classical foundations.

Let \(H\) be the Hilbert class field of \(E\), and let
\(F=E^{\bar{\ }}\).  Every \(F\)-automorphism of an algebraic
closure sends \(H\) to the maximal unramified abelian extension
of its image of \(E\), which is again \(E\).  Hence \(H/F\)
is Galois.  Let \(\sigma\in\operatorname{Gal}(H/E)\) be the
Artin image of \([\mathfrak c]^{-1}\).  Apply Chebotarev to
its conjugacy class in \(\operatorname{Gal}(H/F)\), avoiding
the finitely many primes below \((N)\).  The restriction of every
element in this class to \(E\) is the identity; the selected prime
of \(F\) therefore splits into two distinct primes in \(E\).
Choose a prime of \(H\) whose Frobenius is \(\sigma\) itself
(conjugate the chosen prime if necessary), and let \(\mathfrak p\)
be its restriction to \(E\).  Its relative residue degree over
\(F\) is one, so its Artin symbol in \(H/E\) is also
\(\sigma\).  Thus \([\mathfrak p]=[\mathfrak c]^{-1}\).
The two distinct conjugate primes are comaximal; avoidance gives the
second equality in (ASG.3), and the class equality gives the last.
The use of \(H\) proves existence of an ideal in the original
\(E\); the resulting lattice is not extended to \(H\).
% END ASG3

% BEGIN ASG4
\paragraph{ASG4. Primitive negative vectors in an integral hyperbolic plane.}
\label{claim:ASG4}
Suppose \(\Lambda\) is unimodular, \(e,f\) form a hyperbolic
pair, and \(\mathfrak c\ne0\) satisfies
\(\mathfrak c e\subset\Lambda\) and
\(h(f,\Lambda)\subset\mathfrak c\).  If an integral trace-one
element exists, then for every positive integer \(T\) there is
\(z\in(Ee+Ef)\cap\Lambda\) such that
\[
 h(z,z)=-T,\qquad
       \gamma z\in\Lambda\ \Longrightarrow\ \gamma\in\mathcal O_E.
\tag{ASG.4}
\]
\begin{proof}
Choose \(\mathfrak p,\alpha\) from ASG3 with \(N=T\).
Choose \(p\in\mathfrak p\), \(q\in\bar{\mathfrak p}\)
with \(p+q=1\).  If \(t+\bar t=1\), put
\(w=pt+\bar q\bar t\).  Then \(w\in\mathfrak p\)
and \(w+\bar w=1\).  Define
\[
             b=\overline{-Tw/\alpha},\qquad z=\alpha e+bf.
\tag{ASG.5}
\]
Since \(\alpha\in\mathfrak c\), the first summand lies in
\(\Lambda\).  For \(y\in\Lambda\),
\(h(bf,y)=-Tw h(f,y)/\alpha\) is integral because
\(w\mathfrak c\subset\mathfrak p\mathfrak c=(\alpha)\).
Unimodularity therefore puts the second summand in \(\Lambda\).
The two cross terms give \(h(z,z)=-T(w+\bar w)=-T\).

Suppose \(\gamma z\in\Lambda\).  Pairing with \(f\)
gives \(\gamma\alpha\in\mathfrak c\), equivalently
\(\gamma\mathfrak p\subset\mathcal O_E\).  In particular
\(\gamma w\) is integral.  Pairing with \(\alpha e\in\Lambda\)
gives \(-T\gamma\bar w\in\mathcal O_E\).
Choose \(a\in\mathcal O_E\), \(p_0\in\mathfrak p\)
with \(aT+p_0=1\).  Then
\[
        \gamma=\gamma w+\bar w(\gamma p_0)
                              +a(T\gamma\bar w)
\]
is integral.  This proves actual primitiveness, not merely a nonzero
reduction at the two distinguished primes.
\end{proof}
% END ASG4

% BEGIN ASG5
\paragraph{ASG5. The extra isotropic direction supplied by the actual Forms0 data.}
\label{claim:ASG5}
The original \(\mathrm{Forms0}(D,D.u_0)\) supplies a nonzero
isotropic vector orthogonal to every original hyperbolic pair in \(D.sp\).
\begin{proof}
Its \(\mathrm{val0.de}\) field gives an invertible matrix
\(C\) over \(E\) with \(C^*C\) equal to the original Gram
matrix on \(D.sp.X(0)\).  That space has positive rank by the
rank inequalities of \(D\).  The inverse image under \(C\)
of one standard unit vector gives \(x\) with \(h(x,x)=1\).
Its \(\mathrm{atT.piU}\) field gives \(\pi\in\mathcal O_E\)
with \(\pi\bar\pi=T_{u_0}\).  Let \(v\) be the image of
\(\pi x\) under the original isometric inclusion into \(V\).
The original vector \(z_m\) is orthogonal to this space and has
norm \(-T_{u_0}\).  Therefore
\(w=v+z_m\) has norm zero.  It is nonzero, since
\(h(w,z_m)=-T_{u_0}\ne0\).  All original hyperbolic vectors
are orthogonal to both summands by the defining \(D.sp\) fields.
This proves the assertion in the original space without an isotropy
or approximation theorem as an additional premise.
\end{proof}
% END ASG5

% BEGIN ASG6
\paragraph{ASG6. A saturated isotropic module with the prescribed ideal factors.}
\label{claim:ASG6}
There are independent isotropic vectors \(e_1,\ldots,e_{s+1}\)
and a nonzero fractional ideal \(\mathfrak c\) for which
\[
 \Lambda\cap U=\bigoplus_{i=1}^{s}\mathfrak a_i e_i
                             \ \oplus\ \mathfrak c e_{s+1},
 \qquad U=\operatorname{span}_E(e_1,\ldots,e_{s+1}).
\tag{ASG.6}
\]
\begin{proof}
Adjoin the vector of ASG5 to the original \(s\) isotropic
vectors.  Pairing a linear relation with their original dual vectors
kills the first \(s\) coefficients, and the last vector is nonzero.
They are independent and their span \(U\) is totally isotropic.
The module \(M=\Lambda\cap U\) is saturated in \(\Lambda\):
if \(ay\in U\), \(0\ne a\in\mathcal O_E\), then \(y\in U\).
It has rank \(s+1\), since denominators in a basis of \(U\)
can be cleared into the full lattice \(\Lambda\).

Finite torsion-free modules over a Dedekind domain are projective and
are classified, at fixed positive rank, by their determinant ideal
class.  In the usual Steinitz normal form they are a direct sum of
free rank-one modules and one fractional ideal.  In particular
\(I\oplus J\simeq\mathcal O_E\oplus IJ\) for invertible
fractional ideals.  Choose the class of \(\mathfrak c\) to be
\([\det M]\prod_i[\mathfrak a_i]^{-1}\).
The normal form gives an isomorphism
\(\bigoplus_i\mathfrak a_i\oplus\mathfrak c\simeq M\).
Extend it to \(E\); the images of the standard coordinate vectors
are the desired \(e_i\).  This is an equality of the actual
embedded module after choosing its basis, as in (ASG.6).
All vectors remain in the same totally isotropic \(U\).
The statement also covers \(s=0\).
\end{proof}
% END ASG6

% BEGIN ASG7
\paragraph{ASG7. Integral dual hyperbolic vectors without division by two.}
\label{claim:ASG7}
Write \(c_i=\mathfrak a_i\) for \(i\le s\) and
\(c_{s+1}=\mathfrak c\).  There are isotropic dual vectors
\(f_1,\ldots,f_{s+1}\) with
\[
       h(e_i,f_j)=\delta_{ij},\qquad h(f_i,\Lambda)\subset c_i.
\tag{ASG.7}
\]
\begin{proof}
The saturated quotient \(\Lambda/M\) is finite torsion-free,
hence projective.  Split its exact sequence to obtain a retraction
\(\rho:\Lambda\to M\).  The \(i\)-th coefficient of
\(\rho(y)\) lies in \(c_i\).  Extend this coefficient map
linearly to \(V\); nondegeneracy represents it uniquely as
\(h(f_i^0,-)\).  Thus
\(h(f_i^0,e_j)=\delta_{ij}\) and
\(h(f_i^0,\Lambda)\subset c_i\).
Put \(B_{ij}=h(f_i^0,f_j^0)\).  With the trace-one element
\(t\) of ASG2 set
\[
 C_{ij}=\begin{cases}-B_{ij}&i<j,\\-tB_{ii}&i=j,\\0&i>j,
              \end{cases}
 \qquad f_j=f_j^0+\sum_i C_{ij}e_i.
\tag{ASG.8}
\]
Since \(B\) is Hermitian,
\(B_{ij}+C_{ij}+\overline{C_{ji}}=0\).
Expanding the pairing proves \(h(f_i,f_j)=0\); pairings with
the \(e_i\) are unchanged.

Here is the integrality check for the correction, including the last
ideal \(\mathfrak c\), which need not be conjugation invariant.
For \(y,l\in\Lambda\),
\(h(e_i,y)\in\bar c_i^{-1}\) because \(c_i e_i\subset\Lambda\).
Consequently
\(h(h(e_i,y)f_i^0,l)=\overline{h(e_i,y)}h(f_i^0,l)\)
is integral.  Unimodularity implies
\(h(e_i,y)f_i^0\in\Lambda\).  Pairing this vector with
\(f_j^0\) gives \(h(e_i,y)B_{ji}\in c_j\).
Each correction \(\overline{C_{ij}}h(e_i,y)\) is this element
times \(-1\), \(-\bar t\), or zero.  Hence it lies in
\(c_j\), proving \(h(f_j,\Lambda)\subset c_j\).
No integral division by two has occurred.
\end{proof}
% END ASG7

% BEGIN ASG8
\paragraph{ASG8. The actual orthogonal lattice splitting and all negative vectors.}
\label{claim:ASG8}
Keep the first \(s\) hyperbolic pairs from ASG7 and put
\(\Lambda_b=\Lambda\cap(\sum_i Ee_i+Ef_i)^\perp\).
Then
\[
 \Lambda=\bigoplus_{i=1}^{s}\mathfrak a_i e_i
              \ \oplus\ \Lambda_b
              \ \oplus\ \bigoplus_{i=1}^{s}\mathfrak a_i^{-1}f_i.
\tag{ASG.9}
\]
For every original word \(u\), there is a primitive
\(z_u\in\Lambda_b\) of norm \(-T_u\).
\begin{proof}
By ASG2 the ideals \(\mathfrak a_i\) are conjugation invariant.
Pairing with \(\mathfrak a_i e_i\) shows
\(h(e_i,\Lambda)\subset\mathfrak a_i^{-1}\).
Conversely, if \(b\in\mathfrak a_i^{-1}\), then
\(h(bf_i,\Lambda)\subset\bar b\mathfrak a_i\subset\mathcal O_E\);
unimodularity gives \(\mathfrak a_i^{-1}f_i\subset\Lambda\).
For \(y\in\Lambda\), subtract
\[
            \sum_i h(f_i,y)e_i+\sum_i h(e_i,y)f_i.
\]
Both sums are integral and the remainder is orthogonal to every
\(e_i,f_i\).  This gives (ASG.9); uniqueness follows by pairing
with the hyperbolic vectors and then taking the orthogonal remainder.

The reserved last pair is orthogonal to all the retained pairs.
It satisfies exactly ASG4 with \(\mathfrak c\), by ASG6--ASG7.
Apply ASG4 for each positive \(T_u\) and choose the resulting
\(z_u\) once for that word.  Every vector is in the same
\(\Lambda\), orthogonal to the retained pairs, and primitive in
\(\Lambda\).  This supplies the entire word-indexed family without
changing the common isotropic flag or any ideal \(\mathfrak a_i\).
\end{proof}
% END ASG8

% BEGIN ASG9
\paragraph{ASG9. Every original FixedHyp field and the full lattice block.}
\label{claim:ASG9}
The preceding data give the complete original
\(\mathrm{UnitaryLatticeData}(D)\), including
\(\mathrm{FixedHyp}\) for every word's flag.
\begin{proof}
Use \((e_i)_{i\le s},(f_i)_{i\le s},\Lambda_b,z_u\) from
ASG8 and let \(F_u\) have precisely that isotropic frame and
negative vector.  Its independence, isotropy, orthogonality and norm
are ASG6--ASG8.  The lattice, unimodularity and membership fields are
the original \(\Lambda\) and ASG8.

The integral intersection with \(U_0=\sum_{i\le s}Ee_i\)
is exactly \(\bigoplus\mathfrak a_i e_i\), either by (ASG.6)
or by the projection in (ASG.9).  At every prime \(P\) choose
\(c_{i,P}\in\mathfrak a_i\) generating
\((\mathfrak a_i)_P\).  Such a global member exists by choosing
a nonzero class in \(\mathfrak a_i/P\mathfrak a_i\); its
local coefficient is a unit.  At primes above \(g\), take
\(c_{i,P}=1\), as ASG2 says that these negative powers of the
different contain one and have unit localization there.
Thus localization of the intersection is generated by
\(c_{i,P}e_i\), with exactly the source's prescribed value one
at every distinguished prime.

The remaining primitive condition is slightly stronger in notation
than primitiveness in \(\Lambda\), and can be checked directly.
If \(\gamma z_u=l+x\), with
\(l\in\Lambda\cap U_0^\perp\) and \(x\in U_0\), subtract
\(\sum_i h(f_i,l)e_i\) from \(l\).  This subtraction is
integral, by ASG7, and the remainder is orthogonal to both frames.
Pairing the displayed equality with each \(f_i\) shows that
\(x+\sum_i h(f_i,l)e_i=0\).  Hence \(\gamma z_u\in\Lambda\),
and ASG8 gives \(\gamma\in\mathcal O_E\).
This is precisely primitiveness in the source's \(Lsub\).
All \(\mathrm{FixedHyp}\) fields have now been supplied.
The remaining fields \(\mathrm{Fu\_e}\) and
\(\mathrm{flagZVector}\) hold by the definitions of \(F_u\).
\end{proof}
% END ASG9

% BEGIN ASG10
\paragraph{ASG10. The finite additive character with its exact infinity normalization.}
\label{claim:ASG10}
For every original \(\mathrm{AutInst}\), and hence for ASG1,
there is the complete original \(\mathrm{AdditiveCharacterData}\).
\begin{proof}
On \(\mathbb Q_p\) let \(\{x\}_p\in\mathbb Z[1/p]/\mathbb Z\)
be its negative-power part.  Define
\[
 \psi_{\mathbb Q,f}((x_p))
       =\exp\!\left(-2\pi i\sum_p\{x_p\}_p\right).
\tag{ASG.10}
\]
Only finitely many terms are nonzero, and their classes are additive.
This is a continuous unitary additive character, trivial on
\(\prod_p\mathbb Z_p\).  For \(x\in\mathbb Q\), subtracting
its finite principal parts leaves an integer, so
\(\psi_{\mathbb Q,f}(x)e^{2\pi ix}=1\).
Identify \(E\otimes_F\mathbb A_{F,f}\) with
\(\mathbb A_{E,f}\) and compose (ASG.10) with the adelic trace
\(\operatorname{Tr}_{E/\mathbb Q}\).
The trace of an integral local element is integral, proving
\(\psi_{E,f}=1\) on the actual \(\mathrm{IntAf}\).

The original infinity factor is
\[
 \prod_v \exp\bigl(2\pi i(\phi_v(x)+\overline{\phi_v(x)})\bigr)
               =\exp(2\pi i\operatorname{Tr}_{E/\mathbb Q}x),
\]
because the \(\phi_v\)'s select exactly one member of each CM
pair.  The rational product identity just proved gives the literal
\(\mathrm{psiEf\_rat}\), with its stated sign and without an
extra factor of two.  Additivity and integral triviality give the other
two fields.  Thus this also proves the entire generic
\texttt{ScPsiExists.Target}.
\end{proof}
% END ASG10

% BEGIN ASG11
\paragraph{ASG11. Split local lattices and the literal congruence predicates.}
\label{claim:ASG11}
At every finite place \(v\) of \(F\) split in \(E\), a
self-dual integral Hermitian lattice is a pair of dual free
\(\mathcal O_{F_v}\)-modules.  For positive integers \(g,m\)
and finite adelic vectors, the original predicate
\(\mathrm{InAt}(g,m)\) means precisely membership
in \(m\widehat\Lambda_v\) at every place dividing \(g\).
\begin{proof}
Write \(R=\mathcal O_{F_v}\).  The integral quadratic algebra
is \(R\times R\), with conjugation interchanging its factors.
The two idempotents decompose the lattice as \(M^+\oplus M^-\).
The Hermitian form is zero within either summand and gives a perfect
pairing between them, by self-duality.  Both summands are finite free
over the discrete valuation ring \(R\).  Thus either choice of a
basis on one side determines a dual basis on the other.

The adelic lattice \(I.Lf\) is the product of the local completed
lattices.  One can verify this from its definition as the integral
adelic span of a finite set of global generators: locally these
generators span, and their finitely many coefficient coordinates can
be chosen independently in the product of local rings.
Consequently \(\mathrm{mL}(m)=m\widehat\Lambda\).
If \(c\) is coprime to \(g\), multiplication by \(c\) is a
unit at every place dividing \(g\), proving one implication in
the asserted interpretation of \(\mathrm{InAt}\).
For the converse, clear the finitely many remaining denominator and
divisibility defects by a positive rational integer \(c\) involving
only rational primes not dividing \(g\).  The definition of a
finite adele makes this a finite operation.  Then
\(c x\in m\widehat\Lambda\) and \(\gcd(c,g)=1\), exactly
the required witness.  Similarly \(\mathrm{ZMod'}(g,a)\)
means \(a\in g\widehat{\mathcal O}_E\), and
for integral vectors and integral supplied generators,
\(\mathrm{InMod}\) is the stated residue span at all places
dividing \(g\).  Integrality supplies the conditions away from
\(g\), where \(g\) is a unit.  These interpretations also
cover \(g=1\).
\end{proof}
% END ASG11

% BEGIN ASG12
\paragraph{ASG12. The full original reference frame by finite adelic construction.}
\label{claim:ASG12}
The generic \texttt{ScRefExists.Target} holds.  In particular,
ASG1 and ASG9 have an actual \(\mathrm{ReferenceFrame}\).
\begin{proof}
Use exactly the hypotheses of that target: \(g\ge1\), splitting
of \(E/F\) at every \(v\mid g\), \(n=2r+k+1\), a
hyperbolic frame \(e_i,f_i\) integral at these places, and the
given self-dual lattice.  At such a place the projection
\(y\mapsto\sum_i h(f_i,y)e_i+h(e_i,y)f_i\) is integral.
It splits off the hyperbolic pairs and leaves a self-dual orthogonal
complement of rank \(k+1\) over \(R\times R\).
Choose bases \(u_0,\ldots,u_k\) and \(v_0,\ldots,v_k\)
of its two summands with
\(h(v_i,u_j)=\delta_{ij}(1,0)\), and put
\[
 e_{0,i}=e_i,\quad f_{0,i}=f_i,\quad z_0=u_k-v_k,
 \quad \beta_i=u_i,\quad\beta_i^d=v_i\ (i<k),
 \quad e_w=(1,0).
\tag{ASG.11}
\]
These are local choices at finitely many places.  At all other places
take all these vectors and \(e_w\) to be zero.  They are genuine
integral finite adeles.  ASG11 proves \(\mathrm{basisApprox}\):
the differences from the rational frames vanish at every place above
\(g\), and their other denominators can be cleared away from \(g\).

At the selected places the Gram equalities for the two frames are
exact.  Also \(h(z_0,z_0)=-1\); every required cross pairing
vanishes; \(e_w^2=e_w\), \(e_w\bar e_w=0\),
\(\beta_i=e_w\beta_i\),
\(\beta_i^d=\bar e_w\beta_i^d\), and
\(h(\beta_i^d,\beta_j)=\delta_{ij}e_w\).
Off these places \(g\) is a unit, so all required congruences
are automatic for integral vectors.  The exact idempotent and support
equalities also hold there because both sides are zero.
The condition \((1-e_w-\bar e_w)y\in g\widehat\Lambda\)
holds place by place: its left side is zero above \(g\) and
the modulus is a unit elsewhere.

Finally the displayed vectors span modulo \(g\).
The first pairs span their hyperbolic part; \(\beta_i,\beta_i^d\)
give the first \(k\) coordinates of the complement; and
\(e_w z_0=u_k\), \(\bar e_w z_0=-v_k\) give its last two
components.  This proves the actual \(\mathrm{latticeSpan}\)
field, including \(k=0\).  Every field of \(\mathrm{RefData}\)
has been verified.  For \(g=1\) the empty local construction gives
zero vectors and zero idempotent, and every congruence is modulo a unit,
so it is an allowed instance as well.
The original \(D\)-specialization follows from ASG2 and ASG9.
\end{proof}
% END ASG12

% BEGIN ASG13
\paragraph{ASG13. Transport to any original reference frame, with exact support.}
\label{claim:ASG13}
The full generic \texttt{ScXuExists.Target} holds for any supplied
reference frame, not just the particular frame of ASG12.
It supplies \(\mathrm{TransportData}\) for every word in ASG9.
\begin{proof}
At a place \(v\mid g\), split off the exact integral hyperbolic
frame as in ASG12 and call its orthogonal complement \(W\).
The given vector \(z\) lies in \(W\) and has norm \(-t\),
where \(t\equiv1\pmod g\).  Reduce modulo \(gR\).
The reference vector has norm \(-1\) and is orthogonal to the
reduced hyperbolic frame, by \(\mathrm{basisApprox}\) and
\(\mathrm{formsAtZero}\).  Its projection to \(W\) therefore
has the same reduction.  In the split description of ASG11 a vector
is a pair \((u,\lambda)\) with \(\lambda(u)=-1\).

Over the local Artinian ring \(R/gR\), general linear matrices act
transitively on such pairs by
\((u,\lambda)\mapsto(Au,\lambda A^{-1})\).
To see this, \(u\) is unimodular and
\((R/gR)^{k+1}=(R/gR)u\oplus\ker\lambda\).
Choose a basis of the free complement; it identifies the pair with
\((e_1,-e_1^*)\).  Do this for both pairs and compose the basis
changes.  Thus choose \(\bar A\) taking the reference pair to
the reduction of \(z\).  Lift its entries to \(R\); its
determinant remains a unit, so the lift is in \(\mathrm{GL}_{k+1}(R)\).
Together with its inverse dual it is an integral unitary operator
on \(W\).  Extend it by the identity on the hyperbolic part.

Use these operators at \(v\mid g\), and the identity elsewhere,
to obtain \(x\in G_f\).  It preserves the entire adelic lattice
and fixes every original \(e_i,f_i\) exactly.  Its inverse sends
\(z\) to the reference \(z_0\) modulo \(g\) at all selected
places, which is the literal \(\mathrm{InAt}(g,g)\) by ASG11.
If \(m\) is coprime to \(g\), then \(x-1\) is zero at
every place above \(m\); elsewhere it preserves integral lattices
and \(m\) is a unit.  Hence \((x-1)y\in m\widehat\Lambda\)
for all \(y\in\widehat\Lambda\), proving \(\mathrm{OneOff}\)
with clearing integer one.  For \(g=1\), use \(x=1\) and
clear the remaining rational denominators as in ASG11.
All four conclusion fields of the target are proved.

For the original specialization, \(t=T_u\) has the required
congruence by the original word datum.  ASG8--ASG9 give its norm,
orthogonality and local integrality.  Choose the resulting \(x_u\)
for every word with one already fixed reference frame.  These are
exactly \(\mathrm{xu\_fix},\mathrm{xu\_z},\mathrm{xu\_one}\)
and \(\mathrm{xu\_K}\).
\end{proof}
% END ASG13

% BEGIN ASG14
\paragraph{ASG14. The actual finite-level character idempotent is admissible.}
\label{claim:ASG14}
For the original two-prime modulus, any reference frame as above, the
character of ASG10, and any positive Haar measure \(\mu_f\), the
literal \(\mathrm{eOf}\) is an integrable
\(\mathrm{IsAdmIdem}(\mu_f,g,g_0)\).
\begin{proof}
Here the exact finite calculation of ARO7--ARO10 uses only the lattice,
reference frame, character and modulus, all now available before any
\(\mathrm{HidaSetting}\).  We recall the argument to make this
independence explicit.  At each conjugate pair of residue places the
nonzero parts of the frames give dual bases, since their pairing matrix
is the identity and \(\mathrm{latticeSpan}\) gives spanning.
Order one basis as \((e;\beta;z;f_{\rm reversed})\); in the
other orientation use the reversed active dual \(\beta\)-frame.
The \(K_2\) congruences make the direct \(e,\beta,z\) columns
upper unitriangular, with the middle entries of the \(z\) column
zero.  Applying the conjugate conditions to the inverse dual matrix
does the same for the remaining columns.  Conversely that matrix
shape gives all defining congruences, at both orientations.

Such matrices form a subgroup.  Every coordinate used in the original
argument
\[
 C(x)=\sum_a h(f_{0,a},x\,\mathrm{next}(e_0,z_0)_a)
                    +\sum_{i<k-1}h(\beta_i^d,x\beta_{i+1})
\tag{ASG.12}
\]
is additive under multiplication, modulo the modulus.  For successive
entries this is ordinary unitriangular matrix multiplication.  For the
link from the last \(e\) to \(z\), the possible intermediate
terms are zero because the intervening entries in the \(z\) column
are zero.  At the conjugate orientation inverse transpose reverses the
chains and changes their signs; these are still additive coordinates.
Thus \(C(xy)-C(x)-C(y)\in g\widehat{\mathcal O}_E\).
It follows that the literal \(\chi(x)=\psi_{E,f}(C(x)/g)\)
is a character of \(K_2\), trivial on \(K_g\).
It has modulus one and factors through the finite residue quotient.

The subgroup \(K_2\) contains \(K_g\) and lies in the compact
open lattice stabilizer.  It is therefore compact open, and
\(v=\mu_f(K_2)\) is positive and finite.  Define exactly
\(e(x)=v^{-1}\overline{\chi(x)}\) on \(K_2\), zero elsewhere.
This is locally constant, compactly supported and integrable.
Convolution on \(K_2\) gives
\(e*e(x)=v\,v^{-2}\overline{\chi(x)}=e(x)\);
outside \(K_2\) no pair of nonzero factors can occur.
Inversion gives \(e(x^{-1})=\overline{e(x)}\), and the kernel
condition gives the prescribed two-sided \(K_g\)-invariance.

For full positive type, group finitely many \(x_i\)'s by their
right cosets \(x_iK_2\).  In a group write \(x_i=x_C k_i\).
Then its contribution to
\(\sum_{i,j}\bar c_i c_j e(x_i^{-1}x_j)\) is
\(v^{-1}|\sum_i c_i\overline{\chi(k_i)}|^2\).
Different groups contribute zero.  This proves nonnegativity for all
finite lists and coefficients, not only positivity of one operator.
Finally every supporting residue matrix is unitriangular, so
\((x-1)^n\) is zero modulo each prime dividing \(g_0\).
This is precisely \(\mathrm{UnipMod}(g_0,x)\).
All seven clauses of \(\mathrm{IsAdmIdem}\) are proved for the
original \(\mathrm{eOf}\), with its actual Haar denominator.
\end{proof}
% END ASG14

% BEGIN ASG15
\paragraph{ASG15. The full automorphic bundle under the original gap.}
\label{claim:ASG15}
For the actual blocks constructed above and the original
\(\mathrm{GapHyp}(D)\), there exists a complete
\(\mathrm{AutomorphicBundle}\).  The separate weighted local
coefficient statement used by one source route is unnecessary for
this bundle's unweighted integrability field.
\begin{proof}
Choose the Haar measure of \(G_f\) with \(\mu_f(K)=1\).
It is sigma finite because the group is second countable and locally
compact.  FV6 supplies a measurable fundamental domain of the actual
rational subgroup; RQC2--RQC6 supply its quotient Hilbert space and
the original strongly continuous unitary right regular representation
\(R\).  These choices fix the measure and fundamental domain for
all fields of this bundle.

At each real place choose the actual discrete-series representation of
DC15 at \(\epsilon_v,\lambda_v=\mathrm{lamOf}(A,v)\).
DC14 verifies its hypotheses from the literal rank bound \(1801\)
and gaps \(16n^{19}\); all casts and sign orders are those of
ASG1.  Choose a nonzero vector in its lowest compact type and normalize
it to \(w_v\).  The coefficient identity of DC11 and decay estimate
of DC6 give
\(g\mapsto\langle w_v,\sigma_v(g)w_v\rangle\in L^1\)
for the original infinity Haar measure.  These statements apply to the
chosen actual representations, so no separate assumption about
\(\mathrm{CoefWeightedLoc}\) is needed.  DC12 supplies the complete
four-vector Schur identity and identifies its positive degree with the
literal \(\mathrm{fdegOf}\).

Take the completed finite Hilbert tensor product of these spaces.
Its elementary tensor map has inner product the product of local inner
products; elementary tensors span densely.  The product representation
acts componentwise and is strongly continuous by density and isometry.
This is the explicit tensor construction also used in RFD4, and proves
every field of \(\mathrm{IsHilbTensor}\).
With \(d_v=\mathrm{fdegOf}(\sigma_v,w_v)\), the exact function is
\[
 \phi_\infty(g)=\left(\prod_v d_v\right)
            \overline{\prod_v\langle w_v,\sigma_v(g_v)w_v\rangle}.
\tag{ASG.13}
\]
Tonelli on the product of absolute values proves its \(L^1\)
integrability, including the degree constants.

ASG14 gives the exact integrable finite test function and every
admissibility field.  For any integrable scalar kernel \(a\) and
unitary representation \(U\), the sesquilinear form
\(\int a(g)\langle x,U(g)y\rangle\,dg\) has absolute value
at most \(\|a\|_1\|x\|\|y\|\).  Riesz representation
therefore gives its unique bounded integral operator, of norm at most
\(\|a\|_1\).  Apply this to the finite restriction of \(R\)
with \(a=e\), and to its infinity restriction with
\(a=\phi_\infty\).  These are the required \(R_e,R_\phi\)
and precisely their weak-integral identities.
We have now supplied \(\mu_f,\mathrm{fund},H_v,\sigma_v,w_v\),
the tensor space and map, \(R\), \(e\)'s admissibility,
\(R_e\), all discrete-series and lowest-vector fields,
\(R_\phi\), and \(\phi_\infty\)'s integrability.
These are all fields of \(\mathrm{AutomorphicBundle}\).
\end{proof}
% END ASG15

% BEGIN ASG16
\paragraph{ASG16. The actual finite representation block on every stabilizer.}
\label{claim:ASG16}
For the lattice block of ASG9 there is the complete original
\(\mathrm{RepresentationData}\), with all original highest weights.
\begin{proof}
At a real place the hyperbolic pairs span a nondegenerate space of
signature \((s,s)\).  Their orthogonal complement has signature
\((k,1)\).  Since \(h(z_u,z_u)=-T_u<0\), the further
orthogonal complement of \(z_u\) is positive definite of dimension
\(k\).  Gram--Schmidt gives an orthonormal basis \(w_{u,v,i}\)
there.  These are exactly the \(\mathrm{wfr\_on}\) and
\(\mathrm{wfr\_perp}\) equations after the original real
coordinate map \(\mathrm{rv}\).

Let \(U\) be the span of the isotropic frame.  The real stabilizer
\(S\) in the source acts on \(U^\perp/U\), fixes the class
of \(z_u\), and acts unitarily on its positive orthogonal
complement.  In the chosen basis this gives a continuous homomorphism
\(q_v:S\to U(k)\).  The subgroup \(N\) of the source is
exactly its kernel: its defining last condition says that the action
on \(U^\perp/U\) is the identity, and the class of \(z_u\)
is already fixed by \(S\).

Write \(\nu_i=D.hcW(u,v).u_i\), with indices starting at zero.
These are the original Harish--Chandra parameters, not the integral
highest weights.  For the all-positive sign sequence,
\(\rho_N=0\) and \(\rho_C(i)=(k-1)/2-i\).  Thus the
highest weight required by the literal \(\mathrm{IsHCTypeU}\) is
\[
       \alpha_i=\nu_i-\frac{k-1}{2}+i.
\]
The original rank \(k\) is even and each \(\nu_i\) is
half-integral, so \(\alpha_i\) is an integer.  The original
interlacing has an integral parameter strictly between
\(\nu_i\) and \(\nu_{i+1}\); hence
\(\nu_i-\nu_{i+1}\ge1\), and
\(\alpha_i-\alpha_{i+1}\ge0\).
Apply the Schur-functor and determinant-power construction of DC1
to this nonincreasing integral sequence \(\alpha\).
The resulting finite-dimensional unitary representation
\(\pi_{u,v}\) of \(U(k)\) has precisely the required
\(\mathrm{lowK}(\mathrm{allPos},\nu)\) highest weight.
It therefore satisfies the full original
\(\mathrm{IsHCTypeU}\) with parameter \(\nu\), including
possibly negative last highest weight through the determinant power.
The finite Hilbert tensor product of their spaces is \(V_u\).
On \(S\) define \(\pi_u(s)=\bigotimes_v\pi_{u,v}(q_v(s))\);
extend the function to all of \(G_\infty\) by zero outside \(S\).
The source imposes its representation and continuity laws only on
\(S\).  There they follow from the homomorphisms \(q_v\),
unitarity, and the finite tensor construction; on \(N\) the
operator is the identity.  This proves every \(\mathrm{IsRepS}\)
clause, including continuity on \(S\).

The subgroup \(H\) fixing all \(e_i,f_i,z_u\) acts only on
the positive complement, so its \(q_v\) matrix is exactly
\(\mathrm{matW}(w_{u,v},v,s)\), by its defining pairings.
Thus elementary tensors satisfy the literal
\(\mathrm{tensU\_equiv}\) formula.  Their inner products are
the products of local inner products and they span the full tensor
space (which is finite dimensional).  This proves both other tensor
fields and completes the original representation block for every word.
\end{proof}
% END ASG16

% BEGIN ASG17
\paragraph{ASG17. The original subgroup and quotient measures.}
\label{claim:ASG17}
The same lattice block has a \(\mathrm{MuBlocks}\) object with
the literal infinity Haar measure in its quotient integral identity.
\begin{proof}
The source's groups \(S,N,H\) are closed real matrix subgroups:
their defining flag inclusions are closed linear conditions and their
unitary conditions are closed equations.  The map \(q\) of ASG16
has kernel \(N\), and its restriction to \(H\) is an isomorphism
onto \(\prod_v U(k)\).  Extending an operator on the positive
complement by the identity on \(e_i,f_i,z_u\) proves surjectivity.
Consequently \(S=N\rtimes H\), with multiplication a homeomorphism
from \(N\times H\) onto \(S\): the inverse sends \(s\) to
\((s q(s)^{-1},q(s))\), using this section.

In a basis adapted to \(U\subset U^\perp\), every element of
\(N\) is upper unitriangular.  Indeed the flag action has diagonal
one on \(U\), its middle action is the identity by definition,
and unitarity gives the inverse dual action on the last block.
Thus \(N\) is a real unipotent group and is unimodular.
For completeness its Lie algebra is nilpotent; every adjoint
endomorphism has trace zero, so the modular character is one on its
identity component.  The exponential and finite logarithm polynomials
identify a real unipotent group with its nilpotent Lie algebra, proving
connectedness and the assertion on the whole group.
The group \(H\) is compact.  Its action on a Haar measure of \(N\)
has a positive scalar modulus, a continuous homomorphism
\(H\to\mathbb R_{>0}\).  A compact subgroup of
\(\mathbb R_{>0}\) is trivial, so this action preserves that Haar
measure.  The pushforward of \(\mu_N\times\mu_H\) is therefore
a left and right Haar measure \(\mu_S\) on \(S\).

Regard \(\mu_N,\mu_H\) as measures on the ambient group by
closed inclusion.  They are sigma finite, nonzero, supported on their
subgroups and invariant under their left translations, which are the
four \(\mathrm{IsHaarOnSet}\) fields.
The classical quotient Haar theorem applies to the unimodular closed
subgroup \(S\) of the unimodular real reductive group \(G_\infty\).
Choose its quotient measure relative to this \(\mu_S\) and the
original \(\mu_\infty\).  The second-countable quotient has a
Borel section; push that quotient measure forward to a measure
\(\nu_Q\) on \(G_\infty\).  Weil's identity is then
\[
 \int_{G_\infty} f\,d\mu_\infty
       =\int_{G_\infty}\int_S f(sx)\,d\mu_S(s)\,d\nu_Q(x)
\tag{ASG.14}
\]
for every nonnegative measurable \(f\).  This standard form follows
first for continuous compactly supported functions and then for all
nonnegative measurable functions by equality of the induced Radon
measures and monotone approximation.  It is precisely
\(\mathrm{IsQuotMeasS}\).  The ambient measure has not been
rescaled: the normalization freedom was used in choosing the quotient
measure.  Choose these measures for each word.  This supplies every
\(\mathrm{MuBlocks}\) field, without yet asserting a global
Siegel--Weil measure comparison.
\end{proof}
% END ASG17

% BEGIN ASG18
\paragraph{ASG18. Exact supplied blocks and the remaining common-setting boundary.}
\label{claim:ASG18}
For one original \(D\), its supplied \(\mathrm{Forms0}(D,D.u_0)\),
and \(\mathrm{GapHyp}(D)\), the preceding proofs supply,
simultaneously and on the same objects, the original blocks
\(A,L,P,R_f,X_u,B_u,\Pi,M_\mu\).
They also supply \(\mathrm{RowsFinDim}\) for its actual bundle
through RFD8.  This does not alone assert the full common
\(\mathrm{ConcreteSettingExists}\) or any entire remaining global leaf.
\begin{proof}
Choose \(A\) by ASG1, then \(L\) by ASG9, then \(P\)
by ASG10.  Choose one \(R_f\) by ASG12 and use it for all the
\(X_u\)'s of ASG13 and the bundle of ASG15.  Choose the
representation and measure blocks by ASG16--ASG17 on that same
\(A,L\).  Each construction retains the preceding choices, so
these are dependent blocks of the original constructor, rather than
unrelated existential witnesses.  All word quantifiers occur after
the common lattice and frame choices.  Apply the universal RFD8 to
the actual measure, fundamental domain, finite kernel and infinity
representations now constructed.  It supplies the exact
finite-dimensionality premise used by ARO18 and subsequent row proofs.

The constructor still requires the original nonzero Bessel-function
block \(W\) with its normalized quotient integral and decay, and
the original measured-field block.  The global endpoint additionally
requires the original Siegel--Weil rescaling comparison and its
same-object measure and character ties.  The spectral endpoint also
uses its Bessel uniqueness and local-tie inputs.  These are retained
as their actual source statements; none is inferred from the algebraic
frame or from the existence of Haar measures.  The independent
\(\mathrm{CountIdempotentsClass}\) producer and a common family
threshold likewise remain outside this chapter.
Thus the exact setting-building inputs listed above are now supplied;
substituting them into the original conditional constructors gives the
claimed reduction of those constructors' outstanding inputs, and no
stronger assertion.
\end{proof}
% END ASG18
